\documentclass{article}

\usepackage[margin = 1in]{geometry}

\usepackage{amsmath}

\title{Ch. 14 Non-Textbook Questions}
\author{Principles of Microeconomics}
\date{\today}

\begin{document}

\maketitle

\begin{enumerate}

\item You produce widgets. Currently you produce four widgets at a total cost of \$40.

	\begin{enumerate}
	
	\item What is your average total cost?
	
	\item Suppose you could produce one more (the fifth) widget at a marginal cost of \$5. If you do produce that fifth widget, what will your average total cost be? Has your average total cost increased or decreased? Why?
	
	\item Suppose instead that you could produce one more (the fifth) widget at a marginal cost of \$20. If you do produce that fifth widget, what will your average total cost be? Has your average total cost increased or decreased? Why?
	
	\end{enumerate}
	
\item Don owns a small concrete-mixing company. His fixed cost is the cost of the concrete-batching machinery and his mixer trucks. His variable cost is the cost of the sand, gravel, and other inputs for producing concrete; the gas and maintenance for the machinery and trucks; and his workers. He is trying to decide how many mixer trucks to purchase. He has estimated the costs shown in the accompanying table based on estimates of the number of orders his company will receive per week.
	\begin{center}
	\begin{tabular}{ c | c | c | c | c }
	& & \multicolumn{3}{c}{\textbf{VC}} \\
	\hline
	\textbf{Quantity of trucks} & \textbf{FC} & \textbf{20 orders} & \textbf{40 orders} & \textbf{60 orders} \\
	\hline
	2 & \$6,000 & \$2,000 & \$5,000 & \$12,000 \\
	\hline
	3 & 7,000 & 1,800 & 3,800 & 10,800 \\
	\hline
	4 & 8,000 & 1,200 & 3,600 & 8,400
	\end{tabular}
	\end{center}
	\begin{enumerate}
	
	\item For each level of fixed cost, calculate Don's total cost for producing 20, 40, and 60 orders per week.
	
	\item If Don is producing 20 orders per week, how many trucks should he purchase and what will his average total cost be? Answer the same questions for 50 and 60 orders per week.
	
	\end{enumerate}
	
	Assume that Don purchased 3 trucks, expecting to produce 40 orders per week.
	
	\begin{enumerate}
	
	\setcounter{enumii}{2}
	\item Suppose that in the short run, business declines to 20 orders per week. What is Don's average total cost in the short run? What will his average total cost per order in the short run be if his business booms to 60 orders per week?
	
	\item What is Don's long-run average total cost for 20 orders per week? Explain why his short-run average total cost of producing 20 orders per week when the number of trucks is fixed at 3 is greater than his long-run average total cost of producing 20 orders per week.
	
	\end{enumerate}
	
(Flip) \newpage
	
\item Renata wants to sell earrings this year. It costs \$1,000 to lease an earring making machine for a year. Running the machine costs \$6 per set of earrings produced. Each set takes half an hour to make. Renata values her time at \$15 per hour for the first 1,500 hours she works and \$24 per hour for any additional hours she works. The market price for a set of earrings is \$15. The interest rate on a savings account is 5\%.

	\begin{enumerate}
	
	\item How much money does Renata forgo by taking money out of the bank to lease the earring making machine?
	
	\item What is her marginal cost of making a set of earrings? How does it depend on the quantity produced?
	
	\item To maximize her profit, Renata should produce until the marginal cost of producing another set is greater than the revenue she would earn from selling that set. How many sets of earrings should she produce?
	
	\end{enumerate}
	
	For the rest of the problem assume that she produces the amount you found in part (c).
	
	\begin{enumerate}
	
	\setcounter{enumii}{3}
	\item What are Renata's total explicit costs?
	
	\item What are her total implicit costs?
	
	\item What is her economic profit?
	
	\item What is her accounting profit?
	
	\end{enumerate}
	
\item Evaluate each of the following statements. If a statement is true, explain why; if it is false, identify the mistake and try to correct it.

	\begin{enumerate}
	
	\item A decreasing marginal product tells us that marginal cost must be rising.
	
	\item An increase in fixed cost increases the efficient scale of the firm.
	
	\item An increase in fixed cost increases marginal cost.
	
	\item When marginal cost is above average total cost, average total cost must be falling.
	
	\end{enumerate}

\end{enumerate}

\end{document}